\documentclass[12pt,a4paper,openany,fontset=none]{ctexrep}

\usepackage[
  left=2.8cm,
  right=2.8cm,
  top=2.6cm,
  bottom=2.6cm
]{geometry}
\usepackage{booktabs}
\usepackage{tabularx}
\usepackage{array}
\usepackage{longtable}
\usepackage{xcolor}
\usepackage{enumitem}
\usepackage{hyperref}
\usepackage{fancyhdr}
\usepackage{titlesec}
\usepackage{setspace}
\usepackage{tikz}
\usepackage{pgfplots}
\usepackage{graphicx}
\usepackage{float}
\usepackage{amsmath}

\setCJKmainfont{PingFang SC}
\setCJKsansfont{PingFang SC}
\setCJKmonofont{PingFang SC}

\pgfplotsset{compat=1.18}

\definecolor{MainOrange}{HTML}{E96A2C}
\definecolor{DarkBrown}{HTML}{4B2E1F}
\definecolor{SoftOrange}{HTML}{FFF3E8}
\definecolor{AccentBlue}{HTML}{2F6FED}
\definecolor{DarkGray}{HTML}{333333}
\definecolor{LightGray}{HTML}{F6F6F6}

\hypersetup{
  colorlinks=true,
  linkcolor=MainOrange,
  urlcolor=AccentBlue,
  citecolor=MainOrange,
  pdfauthor={OrangeVC 项目团队},
  pdftitle={OrangeVC AI 创意服务交易平台商业计划书}
}

\titleformat{\chapter}
  {\Huge\bfseries\color{MainOrange}}
  {\thechapter}
  {1em}
  {}

\titleformat{\section}
  {\Large\bfseries\color{DarkBrown}}
  {\thesection}
  {0.8em}
  {}

\titleformat{\subsection}
  {\large\bfseries\color{DarkGray}}
  {\thesubsection}
  {0.8em}
  {}

\pagestyle{fancy}
\setlength{\headheight}{15pt}
\fancyhf{}
\fancyhead[L]{OrangeVC 商业计划书}
\fancyhead[R]{\leftmark}
\fancyfoot[C]{\thepage}

\setstretch{1.25}
\setlist[itemize]{itemsep=0.25em, topsep=0.4em}
\setlist[enumerate]{itemsep=0.25em, topsep=0.4em}

\newcolumntype{Y}{>{\raggedright\arraybackslash}X}
\newcommand{\source}[1]{\par{\small\color{DarkGray}依据：\texttt{#1}}\par}

\begin{document}

\begin{titlepage}
  \centering
  \vspace*{2.5cm}
  {\Huge\bfseries\color{MainOrange} OrangeVC\par}
  \vspace{0.6cm}
  {\LARGE AI 创意服务交易平台商业计划书\par}
  \vspace{1.2cm}
  {\large 基于 GitHub 仓库 \texttt{Orange-ui0811/BigOrange} 当前 \texttt{main} 分支内容生成\par}
  \vspace{0.5cm}
  {\large 版本日期：2026 年 7 月 12 日\par}
  \vfill
  \begin{tikzpicture}
    \fill[SoftOrange] (0,0) circle (2.1);
    \fill[MainOrange] (0,0) circle (1.65);
    \draw[white,line width=2pt] (-0.9,0.7) -- (0.9,0.7);
    \draw[white,line width=2pt] (-1.05,0.1) -- (1.05,0.1);
    \draw[white,line width=2pt] (-0.75,-0.5) -- (0.75,-0.5);
    \fill[DarkBrown] (-0.45,0.25) circle (0.08);
    \fill[DarkBrown] (0.45,0.25) circle (0.08);
    \draw[DarkBrown,line width=1.3pt] (-0.35,-0.25) .. controls (0,-0.55) .. (0.35,-0.25);
  \end{tikzpicture}
  \vfill
  {\small 本文只使用仓库中可核验的信息；仓库未提供的财务预测、真实用户规模、收入流水和融资信息未作虚构。}
\end{titlepage}

\tableofcontents
\clearpage

\chapter{执行摘要}

OrangeVC 是一个面向 AI 创意开发者、AI 实习生、自由开发者与项目需求方的创意服务交易平台。仓库 README 将其定义为“AI 创意服务交易平台”，系统架构文档进一步描述为“Vibe Coding 作品交易与职业匹配平台”。平台以作品展示、需求发布、沟通协商、报价下单、交付验收为交易主线，并叠加 GitHub 榜单、RAG 知识库、趋势分析、方案设计、岗位匹配、AI 自动审核、配额消费等 AI 增强能力。

\source{README.md；docs/OrangeJD系统架构设计说明书.md}

从当前仓库实现看，OrangeVC 已不是概念原型，而是具备双前端、后端服务、数据库结构、部署文档和线上访问记录的全栈系统：

\begin{itemize}
  \item 用户端：\texttt{orange-web}，提供市场、AI 工具、榜单、趋势分析、方案设计、聊天、发布、个人中心、配额等页面。
  \item 管理端：\texttt{orange-admin}，提供控制台、审核、内容、用户、订单、知识库、岗位审核、AI 配置、定价、用量统计等页面。
  \item 后端：\texttt{orange-server}，基于 Spring Boot 3.3.5、Java 21、MyBatis、Redis Stack、Spring AI、WebSocket 实现。
  \item 数据库：\texttt{schema-v1.sql} 当前声明 22 张业务表，覆盖账号、内容、互动、交易、知识库、榜单、趋势、定价等模块。
  \item API：当前 Controller 目录下有 41 个 Controller 文件，显式 \texttt{@GetMapping}/\texttt{@PostMapping}/\texttt{@PutMapping}/\texttt{@DeleteMapping} 映射共 166 个。
\end{itemize}

\source{orange-web/src/router/index.js；orange-admin/src/router/index.js；orange-server/src/main/java/com/orangejd/server/controller；orange-server/src/main/resources/sql/schema-v1.sql}

本商业计划书的重点不是虚构融资模型，而是基于仓库现有产品能力，说明 OrangeVC 的商业定位、用户场景、可交付产品范围、技术壁垒、运营闭环和后续商业化路径。

\chapter{项目定位与用户问题}

\section{项目定位}

OrangeVC 的核心定位是“AI 创意服务交易平台”。它不是单纯的代码托管平台，也不是传统外包平台，而是把 AI 开发作品、项目需求、沟通交易、AI 趋势洞察和岗位能力参考放在同一平台内。

平台面向四类角色：

\begin{longtable}{p{3.2cm} p{10.5cm}}
\toprule
角色 & 核心诉求 \\
\midrule
AI 创意开发者 & 展示作品、沉淀个人主页、获得收藏评论、与需求方沟通、通过报价和订单完成交易。 \\
项目需求方 & 发布需求、浏览作品、联系作者、接受报价、按阶段推进交付验收。 \\
AI 实习生/学习者 & 查看 GitHub 热门项目、趋势图谱、岗位方向和能力要求，用 AI 方案设计拆解学习路径。 \\
平台管理员 & 审核内容、管理用户和订单、配置 AI 审核/知识库/定价/充值码，维护平台内容质量和交易秩序。 \\
\bottomrule
\end{longtable}

\source{README.md；docs/OrangeJD系统架构设计说明书.md；orange-web/src/router/index.js；orange-admin/src/router/index.js}

\section{待解决问题}

仓库架构文档给出了四个明确痛点：

\begin{enumerate}
  \item AI 开发者缺少面向商业价值的作品展示渠道。
  \item 需求方在传统外包平台中难以按 AI 技能、项目类型、作品能力精准找人。
  \item GitHub Trending、招聘 JD、技术文章分散，技术趋势信息割裂。
  \item AI 实习生缺少面向真实岗位的技能参照和学习路径。
\end{enumerate}

OrangeVC 的产品解法是：以交易平台承接供需，以榜单和知识库承接趋势，以方案设计和岗位匹配承接学习成长，以管理后台和 AI 审核承接平台治理。

\source{docs/OrangeJD系统架构设计说明书.md}

\chapter{产品能力与业务闭环}

\section{用户端功能}

当前用户端路由显示，\texttt{orange-web} 已覆盖如下页面：

\begin{longtable}{p{3.5cm} p{10.2cm}}
\toprule
页面/模块 & 功能说明 \\
\midrule
市场 & 浏览作品与需求，是内容发现和交易入口。 \\
AI 工具 & 汇总趋势分析、方案设计、榜单等 AI 能力入口。 \\
榜单 & 展示本平台作品/创作者榜，以及 GitHub 仓库榜。 \\
趋势分析 & 展示 Topic 图谱、AI 分析结果和历史记录。 \\
方案设计 & 将项目想法拆成需求理解、成功标准、风险、任务和岗位匹配等阶段。 \\
作品/需求详情 & 查看详情、评论、联系作者、进入交易或沟通流程。 \\
聊天 & 支持会话、消息、未读提醒和与报价流程衔接。 \\
发布/编辑 & 支持作品和需求发布，进入审核状态流转。 \\
个人中心 & 管理作品、订单、收藏、历史、通知、配额、资料和设置。 \\
岗位投稿 & 用户提交岗位信息，经管理端审核后进入岗位能力体系。 \\
\bottomrule
\end{longtable}

\source{orange-web/src/router/index.js}

\section{管理端功能}

当前管理端路由显示，\texttt{orange-admin} 已覆盖平台运营所需的主要后台能力：

\begin{longtable}{p{3.5cm} p{10.2cm}}
\toprule
页面/模块 & 功能说明 \\
\midrule
控制台 & 查看平台概览，并配置交易模式、AI 托管等关键开关。 \\
审核管理 & 处理作品和需求审核，支持 AI 自动审核状态查看、触发和人工覆盖。 \\
内容管理 & 查看作品与需求详情，执行下架等治理动作。 \\
用户管理 & 用户列表与账号状态控制。 \\
订单管理 & 查看订单详情和订单列表。 \\
知识库 & 管理 RAG 文档，查看统计、上传、列表、删除和详情。 \\
充值码管理 & 生成和查看配额充值码。 \\
岗位审核 & 审核用户投稿岗位，支持导入、同步知识库、通过、拒绝和删除。 \\
AI 调用设置 & 保存 AI 调用配置。 \\
AI 消费定价 & 管理 AI 服务定价，并记录价格变更。 \\
用量统计 & 查看配额使用统计。 \\
\bottomrule
\end{longtable}

\source{orange-admin/src/router/index.js；orange-server/src/main/java/com/orangejd/server/controller/admin}

\section{交易闭环}

OrangeVC 的交易链路由内容、沟通、报价、订单、支付阶段和交付记录共同构成。数据库中对应表包括 \texttt{work}、\texttt{requirement}、\texttt{conversation}、\texttt{message}、\texttt{quote}、\texttt{order}、\texttt{payment\_stage}、\texttt{delivery}。

根据项目进度文档，交易有两条路径：

\begin{enumerate}
  \item 协商路径：聊天沟通 $\rightarrow$ 报价 $\rightarrow$ 接受报价 $\rightarrow$ 生成订单。
  \item 直接购买路径：作品详情 $\rightarrow$ 下单 $\rightarrow$ 支付阶段 $\rightarrow$ 交付验收。
\end{enumerate}

订单状态机覆盖 \texttt{CREATED}、\texttt{WAITING\_DEPOSIT}、\texttt{DEPOSIT\_PAID}、\texttt{IN\_PROGRESS}、\texttt{WAITING\_FINAL\_PAYMENT}、\texttt{DELIVERED}、\texttt{FINAL\_PAID}、\texttt{COMPLETED}、\texttt{CANCELLED}、\texttt{DISPUTED} 等阶段。

\source{docs/项目进度记录.md；orange-server/src/main/resources/sql/schema-v1.sql；orange-server/src/main/java/com/orangejd/server/controller/app/OrderController.java；orange-server/src/main/java/com/orangejd/server/controller/app/QuoteController.java}

\section{内容审核与平台治理}

平台内容治理由人工审核和 AI 自动审核共同完成。作品和需求使用状态机：

\[
\mathrm{DRAFT} \rightarrow \mathrm{PENDING\_REVIEW} \rightarrow \mathrm{APPROVED} / \mathrm{REJECTED}
\]

管理端提供作品、需求审核接口，并提供 AI 自动审核开关、测试、立即运行等接口。开发日志显示，AI 审核使用 DeepSeek V4 Pro，覆盖标题规范、描述完整性、违规内容、图片合规等维度；图片审核失败时降级为纯文本审核，API 异常时保留待审状态。

\source{docs/项目进度记录.md；orange-server/src/main/java/com/orangejd/server/controller/admin/AdminAuditController.java}

\chapter{AI 增强能力}

\section{GitHub 榜单与趋势分析}

OrangeVC 将 GitHub 生态作为创意技术趋势来源。当前分类表插入了 15 个作品类 Topic，其中包括 14 个技术 Topic 和 Other：Creative Coding、Game Development、WebGL、Canvas、Web Audio、AI Agents、RAG、Self-Hosted、Home Assistant、Awesome Self-Host、ESP32、Raspberry Pi、Home Automation、3D Printing、Other。

榜单功能支持：

\begin{itemize}
  \item 本平台作品榜：按收藏数等指标排序。
  \item 本平台创作者榜：按创作者表现排序。
  \item GitHub 仓库榜：按 topic 与 period 拉取 GitHub 数据。
  \item 三段式 period：\texttt{all} 全时段、\texttt{week} 本周涨星、\texttt{new} 本周新增。
\end{itemize}

\source{orange-server/src/main/resources/sql/schema-v1.sql；docs/开发日志/2026-07-09-多模块增强.md；orange-server/src/main/java/com/orangejd/server/controller/app/AppLeaderboardController.java}

\section{RAG 知识库}

仓库集成 Spring AI 1.0.0、Redis Stack Vector Store 和自定义 \texttt{GiteeAiEmbeddingModel}。知识库能力包括文档上传、列表、删除、统计、详情查看，以及 GitHub README 自动采集和向量化。

README 和开发日志显示，系统使用 \texttt{bge-small-zh-v1.5} 作为中文 embedding 模型；为解决中文问题与英文 README 的跨语言检索效果问题，项目采用“入库时将 README 翻译为中文”的策略，并区分全时段永久文档、本周/新增 30 天生命周期文档。

\source{README.md；docs/开发日志/2026-06-24-RAG集成.md；docs/开发日志/2026-06-27-RAG优化.md；orange-server/src/main/java/com/orangejd/server/ai/GiteeAiEmbeddingModel.java；orange-server/src/main/java/com/orangejd/server/controller/admin/AdminKnowledgeController.java}

\section{方案设计与岗位匹配}

用户端 \texttt{PlanDesigner.vue} 与后端 \texttt{AppPlanController} 提供方案设计能力。根据项目文档，该功能按六阶段生成项目规划：想法、需求理解、成功标准、约束风险、任务计划、岗位匹配。岗位模块提供方向列表、岗位列表、岗位详情、用户投稿和 AI 匹配接口；管理端提供岗位导入、审核、同步知识库、编辑和删除。

该能力的商业价值在于把平台从“交易撮合”延伸到“能力成长与项目规划”，使实习生、创作者和需求方都能获得更结构化的决策辅助。

\source{README.md；orange-web/src/views/PlanDesigner.vue；orange-server/src/main/java/com/orangejd/server/controller/app/AppPlanController.java；orange-server/src/main/java/com/orangejd/server/controller/app/JobController.java；orange-server/src/main/java/com/orangejd/server/controller/admin/AdminJobController.java}

\section{配额、定价与充值}

当前仓库包含 \texttt{QuotaService}、\texttt{AppQuotaController}、\texttt{AdminPricingController}、\texttt{AdminRedemptionController}、\texttt{EpayController}，并在数据库中包含 \texttt{ai\_service\_pricing}、\texttt{pricing\_change\_log} 等表。用户端有“我的配额”页面，管理端有“AI 消费定价”和“充值码管理”页面。

这说明平台已具备 AI 服务计费基础设施：按服务配置价格、记录价格变更、用户查询余额、兑换充值码、发起在线充值、查看消耗日志。仓库没有提供真实营收或付费用户数据，因此本文不写收入流水和财务预测。

\source{orange-web/src/router/index.js；orange-admin/src/router/index.js；orange-server/src/main/java/com/orangejd/server/controller/app/AppQuotaController.java；orange-server/src/main/java/com/orangejd/server/controller/admin/AdminPricingController.java；orange-server/src/main/java/com/orangejd/server/controller/admin/AdminRedemptionController.java；orange-server/src/main/resources/sql/schema-v1.sql}

\chapter{技术实现与数据资产}

\section{系统架构}

项目采用单体 Spring Boot 后端加双 Vue 前端结构：

\begin{itemize}
  \item \texttt{orange-common}：统一响应、分页结果、JWT 工具、上下文、异常等公共能力。
  \item \texttt{orange-pojo}：Entity、DTO、VO 等数据对象。
  \item \texttt{orange-server}：Controller、Service、Mapper、Task、Config、AI、WebSocket 等核心后端实现。
  \item \texttt{orange-web}：用户端 Vue 3 应用。
  \item \texttt{orange-admin}：管理端 Vue 3 应用。
\end{itemize}

\begin{figure}[H]
\centering
\begin{tikzpicture}[
  node/.style={draw, rounded corners, fill=SoftOrange, minimum width=3.2cm, minimum height=0.9cm, align=center},
  arrow/.style={->, thick, color=MainOrange}
]
\node[node] (web) at (0,3) {用户端\\orange-web};
\node[node] (admin) at (5,3) {管理端\\orange-admin};
\node[node] (nginx) at (2.5,1.6) {Nginx\\静态资源/API/WS};
\node[node] (server) at (2.5,0.2) {Spring Boot\\orange-server};
\node[node] (mysql) at (-1.2,-1.4) {MySQL\\业务数据};
\node[node] (redis) at (2.5,-1.4) {Redis Stack\\缓存/向量};
\node[node] (ext) at (6.2,-1.4) {外部服务\\Gitee AI/GitHub/飞书/Epay/SMTP};
\draw[arrow] (web) -- (nginx);
\draw[arrow] (admin) -- (nginx);
\draw[arrow] (nginx) -- (server);
\draw[arrow] (server) -- (mysql);
\draw[arrow] (server) -- (redis);
\draw[arrow] (server) -- (ext);
\end{tikzpicture}
\caption{OrangeVC 当前仓库对应的系统结构}
\end{figure}

\source{README.md；pom.xml；orange-server/pom.xml；nginx.conf}

\section{技术栈}

\begin{longtable}{p{3.5cm} p{10.2cm}}
\toprule
层级 & 技术 \\
\midrule
后端框架 & Spring Boot 3.3.5，Java 21。 \\
数据访问 & MyBatis 3.0.3，PageHelper 2.0.0，MySQL Connector/J 8.0.33。 \\
缓存与向量 & Spring Data Redis，Jedis，Redis Stack Vector Store。 \\
AI 框架 & Spring AI 1.0.0，自定义 Gitee AI EmbeddingModel。 \\
鉴权安全 & JWT，BCrypt，双通道 Header：\texttt{Admin-Token} 与 \texttt{App-Token}，Cloudflare Turnstile。 \\
实时通信 & Spring WebSocket，用于消息和通知实时推送。 \\
用户端前端 & Vue 3.5.34，Element Plus 2.14.1，Vite 8.0.12，ECharts 6.1.0，Marked。 \\
管理端前端 & Vue 3.5.34，Element Plus 2.14.1，Vite 8.0.12。 \\
外部集成 & GitHub API，Gitee AI，飞书 Webhook，易支付，QQ 邮箱 SMTP。 \\
\bottomrule
\end{longtable}

\source{README.md；pom.xml；orange-server/pom.xml；orange-web/package.json；orange-admin/package.json；orange-server/src/main/resources/application-prod.yml}

\section{数据表与资产沉淀}

当前 \texttt{schema-v1.sql} 声明 22 张表：

\begin{longtable}{p{4cm} p{9.7cm}}
\toprule
数据域 & 表 \\
\midrule
账号与资料 & \texttt{account}，\texttt{user\_profile}。 \\
内容与审核 & \texttt{category}，\texttt{work}，\texttt{work\_sale\_option}，\texttt{requirement}，\texttt{audit\_record}。 \\
互动与沟通 & \texttt{favorite}，\texttt{comment}，\texttt{browse\_history}，\texttt{conversation}，\texttt{message}，\texttt{notification}。 \\
交易与交付 & \texttt{quote}，\texttt{order}，\texttt{payment\_stage}，\texttt{delivery}。 \\
趋势与知识库 & \texttt{leaderboard\_snapshot}，\texttt{trend\_history}，\texttt{knowledge\_document}。 \\
AI 定价与审计 & \texttt{ai\_service\_pricing}，\texttt{pricing\_change\_log}。 \\
\bottomrule
\end{longtable}

这些表使平台能够积累四类关键资产：创意作品资产、需求与交易资产、技术趋势资产、岗位与知识库资产。对商业化而言，这比单纯论坛或展示站更有价值，因为数据可以反哺推荐、审核、定价、技能匹配和平台运营。

\source{orange-server/src/main/resources/sql/schema-v1.sql}

\chapter{商业模式}

\section{当前可支撑的商业化路径}

基于当前代码，OrangeVC 可以支撑三条商业化路径：

\begin{enumerate}
  \item 交易撮合服务：围绕作品售卖、需求报价、订单阶段和交付验收形成平台交易闭环。仓库已实现报价、订单、支付阶段和交付表结构及 API。
  \item AI 增值服务：趋势分析、RAG 知识库、方案设计、岗位匹配等功能通过配额体系消费。仓库已实现配额余额、兑换、定价、用量统计和 Epay 接口。
  \item B 端/教育场景服务：面向 AI 实习生培养、项目训练营或团队技术选型，提供榜单、知识库、岗位能力图谱和方案设计工具。
\end{enumerate}

\section{暂不写入的商业数字}

仓库中没有真实 GMV、付费用户数、月活用户数、获客成本、转化率、毛利率、融资计划等数据。因此本文不提供这些数字。后续如果需要形成对外融资版 BP，应由项目方补充：

\begin{itemize}
  \item 当前注册用户、活跃用户、发布作品数、需求数、订单数。
  \item 实际付费记录、充值码兑换记录、AI 配额消耗记录。
  \item 服务器成本、AI 调用成本、支付通道成本。
  \item 目标客户访谈、真实成交案例、复购或留存数据。
\end{itemize}

这些信息补齐后，才能计算交易佣金、AI 配额收入、订阅收入和企业版服务收入的财务模型。

\chapter{竞争优势}

\section{产品组合优势}

OrangeVC 的差异点不在单个功能，而在“交易 + AI 趋势 + 岗位能力 + 知识库”的组合：

\begin{itemize}
  \item 相比代码托管平台，OrangeVC 更强调作品的商业展示、沟通和交易。
  \item 相比传统外包平台，OrangeVC 增加 AI 技术标签、GitHub 趋势和岗位能力参考。
  \item 相比单一 AI 工具站，OrangeVC 具备用户、内容、订单、审核、知识库和配额的完整平台基础。
  \item 相比普通学习平台，OrangeVC 将学习路径与真实作品、需求和岗位能力连接起来。
\end{itemize}

\section{技术壁垒}

仓库已实现或明确包含以下技术要素：

\begin{itemize}
  \item 双端前端和完整后台，使平台具备运营治理能力。
  \item 166 个显式 REST 映射和 WebSocket 通信，覆盖主要业务链路。
  \item Redis Stack 向量检索和 Spring AI 集成，支撑 RAG 能力。
  \item GitHub Topic 榜单与 README 入库，形成可持续更新的技术趋势知识源。
  \item AI 审核、飞书通知、Turnstile、人机验证、JWT 双通道、Token 黑名单等安全与治理机制。
  \item 配额、定价、充值码、支付回调，为 AI 功能商业化提供基础设施。
\end{itemize}

\source{README.md；docs/开发日志；orange-server/src/main/java；orange-server/src/main/resources/sql/schema-v1.sql}

\chapter{部署状态与交付成熟度}

\section{线上部署记录}

README 给出在线地址：

\begin{itemize}
  \item 用户端：\url{http://47.116.8.45}
  \item 管理后台：\url{http://47.116.8.45/admin}
\end{itemize}

开发日志记录服务器为阿里云 ECS 2 核 4G，系统为 CentOS 7.9；2026 年 7 月 11 日部署环境包含 JDK 21.0.2、MySQL 8.0.45、Redis Stack Docker、Nginx、Spring Boot 3.3.5 jar、用户端和管理端静态文件。

\source{README.md；docs/开发日志/2026-07-11-服务器部署与引导教程.md；docs/服务器部署完整指南.md}

\section{交付成熟度}

从仓库文档看，项目已经经历了账号体系、分类、内容发布、内容发现、聊天通知、交易闭环、管理端、文件上传、WebSocket、安全加固、部署上线、榜单、趋势分析、AI 审核、RAG、岗位匹配、配额和引导教程等多个阶段。

项目进度记录中给出的全链路验收为：

\begin{quote}
注册 $\rightarrow$ 登录 $\rightarrow$ 发布作品 $\rightarrow$ 审核通过 $\rightarrow$ 市场可见 $\rightarrow$ 聊天（WebSocket）$\rightarrow$ 报价 $\rightarrow$ 下单 $\rightarrow$ 付定金 $\rightarrow$ 交付 $\rightarrow$ 付尾款 $\rightarrow$ 订单完成 $\rightarrow$ 管理后台统计。
\end{quote}

\source{docs/项目进度记录.md}

\chapter{风险与下一步计划}

\section{已识别风险}

仓库文档和配置显示，项目仍有以下需要治理的风险：

\begin{longtable}{p{3.2cm} p{10.5cm}}
\toprule
风险 & 说明与处理方向 \\
\midrule
敏感配置治理 & 开发配置中出现过具体邮箱和 AI Key；生产配置已通过环境变量注入。下一步应清理历史敏感信息并轮换密钥。 \\
数据同步风险 & 部署日志明确指出全量导入数据库会覆盖服务器用户数据，日常更新应只执行增量 schema。 \\
移动端适配 & 2026-07-11 日志记录引导教程仅在 PC 端测试，移动端尚未验证。 \\
AI 成本控制 & 已有定价与配额基础，但缺少真实调用成本、失败率和用户付费数据，商业化前需建立成本报表。 \\
GitHub/API 限流 & 榜单依赖 GitHub API，文档中已有间隔调用和本地缓存策略，后续仍需监控失败重试与限流。 \\
交易合规 & 已实现订单和交付状态，但平台规则、退款争议、内容版权和发票税务规则需要业务制度补齐。 \\
\bottomrule
\end{longtable}

\source{docs/开发日志/2026-07-11-服务器部署与引导教程.md；docs/服务器部署完整指南.md；orange-server/src/main/resources/application-dev.yml；orange-server/src/main/resources/application-prod.yml}

\section{建议路线图}

\begin{enumerate}
  \item 短期：清理敏感配置、补充自动化测试、确认线上环境健康检查、完善移动端适配。
  \item 中期：沉淀真实运营指标，包括注册、活跃、发布、审核通过率、订单转化、AI 消耗、充值和支付成功率。
  \item 中期：完善交易规则，包括退款、争议、版权、违规内容、服务交付标准。
  \item 长期：基于真实数据训练推荐与定价策略，形成作品推荐、岗位推荐、AI 成本优化和企业版能力。
\end{enumerate}

\chapter{结论}

基于当前 GitHub \texttt{main} 分支克隆内容，OrangeVC 已具备一个 AI 创意服务交易平台的核心产品形态：双前端、完整后端、22 张业务表、166 个显式 API 映射、RAG 知识库、GitHub 榜单、趋势分析、方案设计、岗位匹配、AI 审核、配额定价和部署文档。

项目当前最适合的商业表达是“以 AI 创意作品交易为入口，以趋势知识库和岗位能力图谱增强供需匹配，以配额和 AI 服务定价承接增值收入”。下一阶段的关键不是继续堆功能，而是补齐真实业务数据、交易规则、成本核算和运营指标，使商业计划从“可交付产品说明”升级为“可验证增长模型”。

\end{document}
